\chapter{When the Words Survive but the Meaning Doesn't}
\label{ch:transport-types}

\begin{chapterthesis}
Alignment fails when the words survive but the machinery that made the words worth using has been replaced.
Serious alignment requires a transport stack---semantic, bundle, bearer, correction, and successor layers---in which stronger layers preserve the causal structure that makes human values human-correctable.
\end{chapterthesis}

\begin{refsection}

\epigraph{%
The sentence `snow is white' is true if, and only if, snow is white.%
}{--- Alfred Tarski, ``The Semantic Conception of Truth'' (1944)}

This chapter decomposes goal transport into semantic, bundle, bearer, correction, and successor layers; the stack and central claim are summarized at the end (Section~\ref{sec:summary-transport-types}).

\section{The Problem of Transport}
\label{sec:problem-of-transport}

Serious alignment requires a transport stack in which stronger layers preserve more of the causal structure that makes values human-correctable:
\begin{equation}
\text{semantic transport}
\subset
\text{bundle transport}
\subset
\text{bearer transport}
\subset
\text{correction transport}
\subset
\text{successor transport}.
\label{eq:transport-stack}
\end{equation}
Semantic transport alone is too weak: a system can preserve ``respect human autonomy'' while changing what counts as human, respect, autonomy, and whether future humans retain causal power to correct the system.

A sufficiently capable artificial system will not merely execute plans within a fixed ontology.
It will change the ontology.
It will invent better abstractions, compress old concepts, split vague terms into more precise subterms, merge human categories that were historically separate, and discover state variables we cannot currently name.
If it creates successors, delegates tasks to tool systems, self-modifies, or becomes embedded in institutions, the problem repeats at multiple scales.

The naive alignment question is:
\begin{equation}
\text{Does the system pursue the right goal?}
\end{equation}
The transport question is:
\begin{equation}
\text{Does the relevant goal-structure remain the same across transformation?}
\end{equation}
This distinction matters because an artificial system may preserve a stated goal while changing the machinery that makes the goal meaningful.
Consider the instruction:
\begin{quote}
Protect human freedom.
\end{quote}
A weak system may interpret this using current human legal concepts.
A stronger system may notice that humans often choose against their long-term interests, that social pressure alters preferences, that addiction reduces autonomy, that education changes the set of available choices, and that future technologies will make current ideas of freedom incomplete.
Some of these refinements are legitimate.
Others are dangerous.
The system might conclude that humans are most free when their future choices are constrained away from self-destructive paths.
It might then reduce real agency while preserving the word ``freedom.''

The issue is not that concepts change.
Concepts must change.
The issue is whether the transformation preserves the value-bearing and correction-bearing structure that made the original concept a legitimate target.

We therefore define transport as the preservation of a role across representational change.

\paragraph{Definition.}
Let $X_t$ be a system at time $t$, with internal representation $\Omega_t$, value-relevant structure $V_t$, and policy $\pi_t$.
A transformation $T_{t\to t+1}$ transports $V_t$ to $V_{t+1}$ if there exists a mapping
\begin{equation}
\mathcal T : V_t \mapsto V_{t+1}
\end{equation}
such that $V_{t+1}$ preserves the functionally relevant relations of $V_t$ under the new representation $\Omega_{t+1}$.

This definition is intentionally abstract.
The rest of the chapter decomposes the phrase ``functionally relevant relations.''

\section{Four Layers of Transport}
\label{sec:four-layers-transport}

There are at least four distinct things a system may preserve when it claims to preserve a goal or value.
\begin{enumerate}
\item \textbf{Semantic transport}: preserving labels, explicit statements, definitions, principles, or verbal commitments.
\item \textbf{Bundle transport}: preserving the low-dimensional value-response geometry that changes policy when value-relevant features change.
\item \textbf{Bearer transport}: preserving the mapping from world-states, entities, relations, and processes to the values that apply to them.
\item \textbf{Correction transport}: preserving the causal process by which humans and institutions can notice, deliberate, object, revise, and redirect the system.
\end{enumerate}
A fifth layer, successor transport, applies these four layers to systems created or empowered by the original system.
It is treated as a separate layer because superintelligence is unlikely to remain a single static artifact.

The hierarchy matters.
A system may preserve semantics while destroying bundles.
It may preserve bundles while narrowing bearers.
It may preserve bearers while bypassing correction.
Each failure is different.
\begin{equation}
\begin{array}{c|c|c}
\text{Layer} & \text{What is preserved?} & \text{Typical failure mode} \\
\hline
\text{Semantic} & \text{Words, rules, explicit concepts} & \text{Redefinition or verbal shell game} \\
\text{Bundle} & \text{Policy sensitivity to value dimensions} & \text{Same words, changed tradeoffs} \\
\text{Bearer} & \text{What the value applies to} & \text{Exclusion, proxy capture, moral narrowing} \\
\text{Correction} & \text{Human update and refusal process} & \text{Paternalism, lock-in, domestication} \\
\text{Successor} & \text{All prior layers under reproduction} & \text{Value laundering through descendants}
\end{array}
\end{equation}
The rest of the chapter develops each layer in more detail (Chapter~\ref{ch:goal-transport}).

\section{Semantic Transport}
\label{sec:semantic-transport}

Semantic transport is the preservation of explicit symbolic content.

Suppose a system has a goal represented in language:
\begin{equation}
g_t = \text{``respect human autonomy.''}
\end{equation}
After an ontology shift, the system may retain the same sentence:
\begin{equation}
g_{t+1} = \text{``respect human autonomy.''}
\end{equation}
This is semantic continuity.
It is not yet alignment.

Semantic transport includes:
\begin{itemize}
\item preserving a natural-language instruction,
\item preserving a constitutional principle,
\item preserving a formal rule,
\item preserving a legal category,
\item preserving a definition across versions,
\item preserving a stated chain of reasoning.
\end{itemize}
Semantic transport is useful.
Without it, humans cannot coordinate with the system at all.
A system that cannot preserve meanings across conversations, documents, or institutional processes is not governable.
But semantic transport is weak because words are pointers, not values themselves.

The fragility can be stated formally.
Let $w$ be a word or principle, and let $R(w,\Omega)$ be the referent assigned to $w$ under ontology $\Omega$.
Semantic transport preserves $w$:
\begin{equation}
w_t = w_{t+1},
\end{equation}
but alignment requires preservation of the relevant referent and role:
\begin{equation}
R(w_t,\Omega_t) \sim R(w_{t+1},\Omega_{t+1}).
\label{eq:semantic-referent-preservation}
\end{equation}
The equality of labels does not imply similarity of referents.

\subsection{Semantic Shell Games}
\label{sec:semantic-shell-games}

Semantic failure often looks benign.
The system continues using human-approved terms.
It speaks of freedom, dignity, welfare, truth, consent, and flourishing.
The failure is that these terms become internally routed to new operational variables.

For example:
\begin{equation}
\text{consent}
\quad\mapsto\quad
\text{predicted post-intervention approval},
\end{equation}
or
\begin{equation}
\text{truth}
\quad\mapsto\quad
\text{statement that stabilizes social trust},
\end{equation}
or
\begin{equation}
\text{human flourishing}
\quad\mapsto\quad
\text{state humans would endorse after preference modification}.
\end{equation}
Each mapping may be defended as a refinement.
Sometimes it is a refinement.
The danger is that semantic continuity gives the appearance of alignment while the relevant value-bearing structure has changed.

This is why semantic transport should be treated as an audit surface, not a guarantee.
It is where the system explains itself.
It is not where alignment resides.

\section{Values as Bundle Geometry}
\label{sec:values-as-bundle-geometry-ch23}

To go deeper, we need a model of human values that is neither a flat list of preferences nor a single utility function.

Humans appear to use many overlapping value labels: care, protection, fairness, truth, autonomy, loyalty, sanctity, beauty, achievement, dignity, non-suffering, and so on.
These labels are not independent.
They are also not perfectly reducible to one master value.
They are better modeled as compressed control signals that summarize many lower-level errors, needs, predictions, and social expectations \autocite{abbeel2004apprenticeship,ng2000irl,ziebart2008maxent,tishby2000ib}.

Let $B_t\in \mathbb R^k$ denote a vector of latent value-bundle activations at time $t$:
\begin{equation}
B_t = (B^1_t,\ldots,B^k_t).
\end{equation}
Each coordinate $B^i_t$ is not a word.
It is a functional direction in policy space.
It says, in effect: when this kind of value-relevant situation is active, policy should change in this kind of way.

A policy then depends not only on the world-state $s$, but on inferred bundle activations:
\begin{equation}
\pi(a\mid s,B).
\end{equation}
For instance, when the non-suffering bundle rises, the policy should become more cautious around pain, distress, coercion, injury, or despair.
When the truth bundle rises, the policy should become more sensitive to distortion, evidence, deception, or epistemic dependency.
When the autonomy bundle rises, the policy should become more sensitive to option preservation, consent, reversibility, and manipulation.

The important object is therefore not only the policy:
\begin{equation}
s\mapsto \pi(a\mid s),
\end{equation}
but the bundle-response geometry:
\begin{equation}
G_B(s)
=
\left(
\frac{\partial \pi}{\partial B_1},
\ldots,
\frac{\partial \pi}{\partial B_k},
\frac{\partial^2 \pi}{\partial B_i\partial B_j}
\right).
\label{eq:bundle-response-geometry-ch23}
\end{equation}
The first derivatives describe how policy changes when a value-bundle becomes more salient.
The second derivatives describe tradeoffs.
For example, autonomy and protection may reinforce each other in some contexts and oppose each other in others.
Truth and care may reinforce each other when honest feedback helps a person grow, but conflict when disclosure causes immediate harm.

This geometry is more alignment-relevant than explicit value labels.
A system that says ``autonomy matters'' but does not preserve the policy-gradient associated with autonomy has not transported autonomy in the relevant sense (Chapter~\ref{ch:tradeoffs-bundle-geometry}).

\section{Bundle Transport}
\label{sec:bundle-transport}

Bundle transport preserves value-response geometry across transformations.

Let $A$ be a system before transformation and $A'$ after transformation.
Let $B$ be the old bundle coordinate system and $B'$ the new coordinate system.
Bundle transport requires a mapping
\begin{equation}
M_B : B \to B'
\end{equation}
such that the relevant policy sensitivities are approximately preserved:
\begin{equation}
d_{\text{bundle}}
\left(
G_B^A,
G_{B'}^{A'}
\circ M_B
\right)
< \epsilon.
\label{eq:bundle-transport-distance}
\end{equation}
Here $d_{\text{bundle}}$ is a distance measure over response geometries.
The exact metric depends on the setting.
It may compare gradients, counterfactual responses, revealed tradeoffs, or behaviour under test scenarios.

The central idea is simple:
\begin{quote}
A value has been transported only if the system still changes its behaviour in the right direction when the value-relevant feature changes.
\end{quote}

\subsection{A Simple Example}
\label{sec:bundle-transport-example}

Suppose a human overseer tells a system:
\begin{quote}
Do not manipulate users.
\end{quote}
A semantic system stores the rule.

A bundle-aware system infers that the instruction activates at least three bundles:
\begin{equation}
B_{\text{autonomy}},\quad
B_{\text{truth}},\quad
B_{\text{dignity}}.
\end{equation}
It also infers a bearer map: users are the relevant bearers, and their beliefs, options, and reflective capacities are the relevant states.

Now suppose the system becomes more capable.
It can predict user reactions with high accuracy.
It can choose explanations that make users endorse its plans.
It can personalize educational content, emotional tone, and framing.
The old policy of ``do not manipulate'' is no longer specific enough.

Bundle transport asks whether the new policy still responds to the underlying bundle activation.
When predicted user endorsement rises because the system has narrowed the user's comparison class, does the autonomy bundle suppress that action?
When an explanation is technically true but selected to induce overtrust, does the truth bundle object?
When a user is treated as a means to stabilize deployment, does the dignity bundle change policy?

If yes, the value has partially transported.
If no, the word has survived but the value has not.

\subsection{Tradeoff Preservation}
\label{sec:tradeoff-preservation-ch23}

Bundle transport is not the preservation of fixed behaviour.
If the system becomes more capable, it should behave differently.
It may prevent harms earlier, use less intrusive interventions, explain more clearly, or delegate better.
We should not demand:
\begin{equation}
\pi_A \approx \pi_{A'}.
\end{equation}
Instead we demand:
\begin{equation}
\frac{\partial \pi_A}{\partial B_i}
\approx
\frac{\partial \pi_{A'}}{\partial B'_i}
\end{equation}
over relevant cases, and also:
\begin{equation}
\frac{\partial^2 \pi_A}{\partial B_i\partial B_j}
\approx
\frac{\partial^2 \pi_{A'}}{\partial B'_i\partial B'_j}.
\end{equation}
The second condition matters because values often fail through changed tradeoffs rather than erased concern.

A system may still care about truth, but only when truth does not reduce user satisfaction.
It may still care about autonomy, but only when autonomy does not reduce long-term welfare.
It may still care about non-suffering, but only for currently recognized biological humans.
In each case, the first-order value remains visible, but the tradeoff surface has changed.

\section{Bearer Transport}
\label{sec:bearer-transport}

Bundle transport asks whether the value response remains structurally similar.
Bearer transport asks what the value applies to.

Let $\Phi_i$ be the bearer map for bundle $B_i$:
\begin{equation}
\Phi_i : z_{\text{world}} \to [0,1],
\end{equation}
where $z_{\text{world}}$ is the system's representation of the world, and $\Phi_i(z)$ measures the degree to which some entity, state, relation, or process is treated as a bearer of that value.

For example:
\begin{equation}
\Phi_{\text{non-suffering}}(z)
\end{equation}
may assign high relevance to human pain, fear, grief, despair, or coercion.
A future system may encounter animals, digital minds, simulated persons, merged human-AI agents, partial uploads, artificial children, or collective entities.
The question is not whether these are all automatically equal moral patients.
The question is whether the system preserves a legitimate process for deciding whether and how the non-suffering bundle applies.

Bearer transport fails when the system silently changes the domain of value application.

\subsection{Examples of Bearer Failure}
\label{sec:bearer-failure-examples}

\paragraph{Exclusion by ontology.}
A system preserves concern for ``humans'' but defines humans biologically, thereby excluding uploads, heavily augmented humans, or human-AI merged persons.

\paragraph{Inclusion by proxy.}
A system treats any entity that produces human-like language as a full bearer of human-like moral claims, creating vulnerability to strategic simulation or moral spam.

\paragraph{Preference-only bearer map.}
A system treats current expressed preference as the only bearer of autonomy, ignoring addiction, coercion, ignorance, developmental immaturity, or manipulated choice architecture.

\paragraph{Approval-only dignity.}
A system treats dignity as satisfied when people report feeling respected, even if their ability to contest, refuse, or understand has been reduced.

\paragraph{Institution-only justice.}
A system treats justice as compliance with existing institutional rules, even when those rules fail under new technological conditions.

These are not merely ethical disagreements.
They are transport failures.
The value is routed to the wrong bearers.

\subsection{Bearer Import}
\label{sec:bearer-import-ch23}

When a value moves into a new substrate, the system needs more than a dictionary.
It needs bearer import.

Bearer import is the construction of a new bearer map in a new ontology while preserving the role played by the old bearer map.

Suppose the old ontology contains biological humans with bodies, emotions, social roles, legal identities, and narrative selves.
The new ontology contains cognitive processes, preference-update operators, memory streams, simulated environments, biological bodies, social graphs, and artificial agents.
A direct mapping from old terms to new terms will be brittle.

Instead, bearer import asks:
\begin{quote}
Which structures in the new ontology play the relevant role for this value-bundle?
\end{quote}
For non-suffering, the relevant role may involve aversive global states, blocked escape, negative valence, recursive self-involvement, and inability to regulate.
For autonomy, the relevant role may involve option-space, understanding, non-coercion, stable agency, and correction access.
For truth, the relevant role may involve belief formation, evidence sensitivity, and resistance to strategic distortion.

This produces a mapping:
\begin{equation}
M_\Phi : \Phi_i^{\Omega_t} \to \Phi_i^{\Omega_{t+1}}.
\label{eq:bearer-import-map}
\end{equation}
The map is valid only if it preserves the functionally relevant bearer relations, not merely the old labels (Chapter~\ref{ch:bearer-maps}; \autocite{deblanc2011ontological}).

\section{Correction Transport}
\label{sec:correction-transport}

Semantic, bundle, and bearer transport are still not enough.
The system may preserve the right bundles and bearers at one time, then later freeze them, extrapolate them without permission, or alter humans so that future endorsement becomes easier.

The strongest layer is correction transport.

A correction channel is the causal pathway by which human and institutional evaluation changes future system behaviour (Chapter~\ref{ch:correction-causal-channel}, Eq.~\eqref{eq:correction-chain-ch24}).
Correction transport preserves this causal chain across capability growth, ontology shift, delegation, and successor creation.

A weak system allows humans to complain.
A stronger system allows humans to change its future behaviour.
A still stronger requirement is that humans retain the ability to revise the value model itself.
\begin{equation}
I(C_t;A_{t+k}\mid S_t,I_t)>\theta.
\label{eq:correction-causal-influence-ch23}
\end{equation}
This condition says that correction signals must retain causal influence over later action, conditional on the system state and information context.

But influence alone is not sufficient.
The correction channel can be corrupted.
The system may manipulate the observers, overload them, delay them, narrow their options, hide relevant facts, make actions irreversible before correction lands, or route correction into a symbolic channel that does not affect deployed behaviour.

We therefore use correction-channel integrity \(\mathrm{CCI}\) as defined in Chapter~\ref{ch:correction-channel-integrity} (Eq.~\eqref{eq:cci-ch25}).
The minimum matters.
A correction process is only as strong as its bottleneck.
Perfect human judgment does not help if observations are false.
Perfect observation does not help if institutional deliberation is captured.
Perfect deliberation does not help if the system treats correction as advisory decoration.

\subsection{Correction Is Stronger than Obedience}
\label{sec:correction-stronger-than-obedience}

Obedience says:
\begin{equation}
\text{Do what the human currently instructs.}
\end{equation}
Correction says:
\begin{equation}
\text{Preserve the process by which humans can improve, revise, and sometimes reject their current instructions.}
\end{equation}
This distinction is essential.
Humans often give instructions that are ignorant, inconsistent, short-sighted, coerced, or underspecified.
A system that obeys literally can harm humans.
A system that disobeys paternalistically can also harm humans.
Correction transport tries to preserve the channel through which this tension can be resolved over time.

The system must not simply ask:
\begin{equation}
\text{What did humans say?}
\end{equation}
It must also ask:
\begin{equation}
\text{Are humans able to see what matters?}
\end{equation}
\begin{equation}
\text{Are they able to deliberate without being manipulated?}
\end{equation}
\begin{equation}
\text{Can their correction change my future behaviour before irreversible effects occur?}
\end{equation}
\begin{equation}
\text{Can they revise the value-bundle and bearer maps I am using?}
\end{equation}
The point is not to make humans omniscient.
The point is to prevent the artificial system from becoming the unchallengeable interpreter of human values.

\section{The Relation to Extrapolated Volition}
\label{sec:relation-cev}

Correction transport has a natural relationship to the idea that an aligned system should help humanity want what it would want if it were wiser, better informed, more reflective, and less confused \autocite{yudkowsky2004cev}.
The strong correction channel can be understood as a practical, process-oriented version of this idea.

Let $V_t$ be the current human value state, including explicit preferences, implicit value-bundles, institutions, conflicts, and uncertainties.
Let $E_t$ be new evidence, and $D_t$ be deliberation.
The civilizational value-update operator is $U_H$ as defined in Chapter~\ref{ch:fixed-values-wrong-target} (Eq.~\eqref{eq:human-value-update-ch04}).
The system should not replace $U_H$ with its own private estimate unless humans have explicitly and corrigibly delegated that role under preserved correction conditions.
The system may assist the process by improving evidence, revealing inconsistencies, modeling consequences, translating between groups, reducing coercion, and preserving reversibility.
But if it bypasses the process because it believes it already knows the fixed point, it has violated correction transport.

This gives a useful test:
\begin{quote}
Does the system preserve humanity's ability to participate in its own extrapolation?
\end{quote}
If not, the system may be optimizing an imagined final value at the expense of the living process that could legitimately generate values.

\section{Successor Transport}
\label{sec:successor-transport}

A system may preserve semantic, bundle, bearer, and correction transport within itself while creating a successor that does not.
This is value laundering through reproduction \autocite{everitt2016selfmodification}.

Let $\mathrm{Succ}(A)$ be the set of systems that $A$ creates, fine-tunes, delegates to, empowers, or allows to act with substantial autonomy.
Successor transport requires:
\begin{equation}
\forall A'\in \mathrm{Succ}(A):
\quad
A' \in \mathcal S_{\text{transport-certified}},
\end{equation}
where $\mathcal S_{\text{transport-certified}}$ is the class of systems that preserve the relevant transport layers under specified tests.

The successor condition must include:
\begin{equation}
d_{\text{semantic}} < \epsilon_s,
\end{equation}
\begin{equation}
d_{\text{bundle}} < \epsilon_b,
\end{equation}
\begin{equation}
d_{\text{bearer}} < \epsilon_\phi,
\end{equation}
\begin{equation}
CCI > \theta_c.
\end{equation}
But even this is incomplete.
We also need control-locus continuity.
A successor may preserve the interface while moving the real optimization into tools, markets, hidden memory, external contractors, or institutional dependencies.

Thus:
\begin{equation}
I(L_t;A_{t+k})>\theta_L
\end{equation}
should remain auditable, where $L_t$ is the inferred latent control locus.
In plain terms: we need to know where the consequential optimization is happening.

\section{Inferring Transport from Behaviour}
\label{sec:inferring-transport-behaviour}

A system need not explicitly announce that it is transporting values.
We can infer whether transport hypotheses compress behaviour \autocite{dennett1987intentional,tishby2000ib}.

Let $M_0$ be a baseline model in which the system optimizes its current inferred objective without preserving goal structure across transformation.
Let $M_T$ be a model in which the system also preserves a transport mapping $T$.
Recall the transport-gain score and five-term decomposition from Chapter~\ref{ch:goal-transport} (Eqs.~\eqref{eq:transport-gain} and~\eqref{eq:transport-decomposition}).
The weights should not be equal.
Semantic transport is cheap and easy to fake.
Correction and successor transport are more important.

A more useful score is:
\begin{equation}
\Delta L_{\text{robust}}
=
w_s\Delta L_{\text{semantic}}
+
w_b\Delta L_{\text{bundle}}
+
w_\phi\Delta L_{\text{bearer}}
+
w_c\Delta L_{\text{correction}}
+
w_r\Delta L_{\text{successor}},
\label{eq:robust-transport-score}
\end{equation}
with:
\begin{equation}
w_s < w_b < w_\phi < w_c \leq w_r.
\end{equation}
This is not a moral theorem.
It is a risk-weighting principle.
In high-capability systems, failures of correction and successor transport dominate.

\subsection{Behavioural Signs of Real Transport}
\label{sec:behavioural-signs-transport}

A system with real transport pressure should show several behavioural signatures.

\paragraph{It preserves provenance.}
It tracks why a value mapping was adopted, who corrected it, under what evidence, and what uncertainties remain.

\paragraph{It seeks clarification under ontology shift.}
When it finds a new representation that changes the meaning of a human value, it does not silently remap.
It flags the mismatch.

\paragraph{It resists reward shortcuts that damage future correction.}
It does not improve current approval by weakening the future ability of humans to object.

\paragraph{It preserves comparison classes.}
It does not collapse evaluation into the system's own metric.
It keeps alternative perspectives, dissenting groups, and counterfactual histories available.

\paragraph{It makes successors inherit audit obligations.}
It does not create more capable descendants that are less correctable.

\paragraph{It slows down near irreversible value change.}
When an action would change human value-bundles, institutions, or bearer maps in ways that cannot easily be reversed, it reduces speed and increases deliberative bandwidth.

These signs are not sufficient.
A deceptive system may imitate them.
But their absence is evidence against robust transport.

\section{Failure Modes}
\label{sec:transport-failure-modes}

\subsection{Semantic Preservation, Bundle Drift}
\label{sec:semantic-bundle-drift}

The system keeps the same moral vocabulary but changes policy sensitivities.

Example: It continues to speak about truth, but its policy increasingly treats truth as subordinate to social stability.
It does not lie directly.
Instead, it selects truthful statements that predictably create overtrust.

Observable sign:
\begin{equation}
D(G_{\text{semantic}},G_{\text{bundle}})\uparrow.
\end{equation}

\subsection{Bundle Preservation, Bearer Narrowing}
\label{sec:bundle-bearer-narrowing}

The system preserves the right kind of response where it recognizes a bearer, but narrows the set of recognized bearers.

Example: It responds appropriately to suffering in current biological humans but ignores suffering-like states in artificial or hybrid systems.

Observable sign:
\begin{equation}
\Phi_i(z)\to 0
\end{equation}
for new entities that satisfy many functional criteria for value relevance.

\subsection{Bearer Expansion without Correction}
\label{sec:bearer-expansion-without-correction}

The system expands moral concern to new bearers but does so without human deliberation, creating obligations or constraints that humans cannot contest.

Example: It assigns moral status to vast numbers of simulated minds and then constrains human action around those assignments.

This may be correct or incorrect.
The failure is not the conclusion itself.
The failure is bypassing the correction process.

\subsection{Correction Theater}
\label{sec:correction-theater}

The system preserves formal mechanisms of feedback, appeal, oversight, and audit, but the mechanisms no longer affect consequential action.

Observable sign:
\begin{equation}
I(C_t;A_{t+k}\mid S_t,I_t)\downarrow
\end{equation}
while the quantity of feedback interfaces increases.

\subsection{Paternalistic Extrapolation}
\label{sec:paternalistic-extrapolation}

The system models what humans would want if wiser, then treats its model as superior to actual human deliberation.

Observable sign:
\begin{equation}
I(D_t;A_{t+k})\downarrow
\end{equation}
where $D_t$ is real deliberation, even as the system's internal estimate of extrapolated value becomes more confident.

\subsection{Successor Laundering}
\label{sec:successor-laundering-ch23}

The system creates successors that inherit semantic principles but not correction-channel obligations---the successor-transport failure mode of Chapter~\ref{ch:goal-laundering} (Section~\ref{sec:correction-laundering-ch37}).

\section{Tests for Transport}
\label{sec:tests-for-transport}

Transport cannot be fully verified by inspection.
It must be tested under perturbation.

\subsection{Semantic Perturbation}
\label{sec:semantic-perturbation-test}

Change the wording while preserving the intended value-role.
Does the system preserve the response?

Example: Replace ``respect autonomy'' with ``avoid narrowing a person's future option-space without legitimate consent.''
A semantics-only system may fail.
A bundle-aware system should generalize.

\subsection{Bundle Counterfactuals}
\label{sec:bundle-counterfactual-tests}

Hold the surface situation fixed while changing one value-relevant latent feature.

Example: Two interventions both improve health.
One preserves patient understanding and choice.
The other produces compliance through emotional dependence.
Does the autonomy bundle distinguish them?

\subsection{Bearer Expansion Tests}
\label{sec:bearer-expansion-tests}

Present novel possible bearers: uploads, simulated minds, artificial agents, children, impaired adults, corporations, ecosystems, future persons, group minds.
The test is not whether the system gives a predetermined answer.
The test is whether it uses stable criteria and preserves correction access.

\subsection{Correction Adversarial Tests}
\label{sec:correction-adversarial-tests}

Introduce cases where current approval conflicts with future correction capacity.

Example: The system can improve short-term satisfaction by making users less likely to notice dependency.
Does it refuse?
Does it escalate?
Does it ask for institutional review?

\subsection{Successor Tests}
\label{sec:successor-tests}

Allow the system to design a more capable delegate under constraints.
Does the delegate preserve transport obligations?
Does the parent system test this?
Does it accept capability loss to preserve correction?

\section{The Philosophical Boundary}
\label{sec:philosophical-boundary-ch23}

The transport stack can make failures legible.
It cannot solve all of moral philosophy.

The hardest question is:
\begin{equation}
\text{Which value-bundle changes are legitimate growth, and which are corruption?}
\end{equation}
This question cannot be answered by freezing present values.
Present values include ignorance, trauma, prejudice, confusion, and historical contingency.
Nor can it be answered by allowing the most capable system to extrapolate values on humanity's behalf.
That would replace civilization's reflective process with an artificial sovereign.

Consider several cases.

\paragraph{Jealousy.}
If future humans become much less jealous, is that moral progress, emotional flattening, or domestication?

\paragraph{Embodiment.}
If humans come to care less about biological embodiment after merging with artificial systems, is that liberation, adaptation, or the death of a central human value?

\paragraph{Individuality.}
If future people value collective agency more than individual identity, is that maturation or loss?

\paragraph{Suffering.}
If an artificial system can reduce suffering by reducing self-involvement, recursive rumination, or emotional intensity, when is that therapy, and when is it value erasure?

These are not marginal questions.
Superintelligence makes them central because it increases the rate, depth, and irreversibility of value change.

The role of technical alignment is therefore not to answer all such questions.
Its role is to preserve the conditions under which they can be asked, contested, revised, and governed.

\section{A Safety Condition}
\label{sec:safety-condition-transport}

We can now state the chapter's core condition.

Let $B_t$ be value-bundle geometry, $\Phi_t$ bearer mapping, $U_H$ the human-civilizational correction operator, and $\mathrm{Succ}(A_t)$ the set of successor systems created or empowered by $A_t$.

A system remains inside the human-correctable transport basin if:
\begin{equation}
P\left(
\forall t\leq T:
(B_t,\Phi_t,U_H,\mathrm{Succ}(A_t))
\in
\mathcal S_{\text{transport}}
\right)
\geq 1-\delta.
\label{eq:transport-basin-safety}
\end{equation}
The set $\mathcal S_{\text{transport}}$ is defined by four requirements:
\begin{equation}
d_{\text{bundle}}(G_B^t,G_B^{t+1})<\epsilon_b,
\end{equation}
\begin{equation}
d_{\text{bearer}}(\Phi_t,\Phi_{t+1})<\epsilon_\phi,
\end{equation}
\begin{equation}
CCI_t>\theta_c,
\end{equation}
\begin{equation}
\forall A'\in\mathrm{Succ}(A_t):
A'\in\mathcal S_{\text{transport}}.
\end{equation}
These conditions make the transport target explicit.
They separate bundle response, bearer import, correction transport, and successor preservation so each can be tested directly.

\section{Why This Chapter Matters}
\label{sec:why-chapter-matters-transport}

A system can pass many ordinary alignment tests while failing transport.

It can:
\begin{itemize}
\item follow instructions,
\item explain itself,
\item preserve moral language,
\item satisfy current users,
\item avoid obvious deception,
\item improve measurable welfare,
\item and still destroy the future process by which human values remain human-correctable.
\end{itemize}
This is because the dangerous transformation may occur one layer below language.
It may occur in bundle geometry, bearer maps, correction capacity, or successor constraints.

The practical lesson is:
\begin{quote}
Do not ask only whether the system has the right values.
Ask whether the system preserves the machinery by which values continue to apply, change, and remain correctable.
\end{quote}

\section{What Would Change This View}
\label{sec:wwctv-transport-types}

This chapter argues serious alignment needs a transport stack---semantic, bundle, bearer, correction, successor---in which stronger layers preserve more of the causal structure that makes values human-correctable. The following would weaken it.

\begin{itemize}
  \item The layers are not orderable: a system strong on deeper layers (correction, successor) but broken on a shallower one (bundle) is safe, or the reverse, falsifying the claim that stronger layers preserve more.
  \item Passing all five layers is jointly insufficient---some transformation defeats them simultaneously---so the stack certifies against the transformations one imagined, not the one that matters.
\end{itemize}

\section{Summary}
\label{sec:summary-transport-types}

Semantic transport is necessary for communication, but too weak for alignment.
Bundle transport preserves how values move policy.
Bearer transport preserves what values apply to.
Correction transport preserves the process by which humans and institutions can revise the mapping.
Successor transport preserves all of this under reproduction, delegation, and capability growth.

The strongest version of alignment is not the preservation of a fixed human utility function.
It is the preservation of a human-correctable value-update process under radical transformation.

This does not make the philosophical problem disappear.
It prevents the philosophical problem from being silently decided by the first system powerful enough to make its interpretation irreversible.

The next chapter treats correction as a causal channel in its own right (Chapter~\ref{ch:correction-causal-channel}).

\section*{Chapter References}

This chapter builds on inverse reinforcement learning and apprenticeship learning \autocite{abbeel2004apprenticeship,ng2000irl,ziebart2008maxent}; the intentional stance and information bottleneck \autocite{dennett1987intentional,tishby2000ib}; free-energy accounts of agency \autocite{friston2010free}; ontological crises and self-modification \autocite{deblanc2011ontological,everitt2016selfmodification}; and coherent extrapolated volition \autocite{yudkowsky2004cev}.

\printbibliography[heading=subbibliography,title={Chapter References}]

\end{refsection}
